\subsection{Common Edits}\label{subsec:common-edits}
This section is going to address a few common changes, that can be applied to the framework.

\subsubsection*{Adding Meta Description to Homepage}
All pages created through administration dashboard include textarea to fill in meta description.
The homepage component is defined in src/components/Home/Home.php.
In the constructor the HtmlHead is created, and it includes function to add meta fields.
Adding description editable through web settings can be added this way:

\begin{verbatim}
    parent::__construct(
        new ContextAwareWebPage(
            // Save the HtmlHead to variable
            head: $head = new HtmlHead(...)
        )
    );

    $description = Setting::fromName(
        'project:homepage-description',
        // Create the setting if it is not defined
        true,
        // Default description
        '',
        // Makes the setting editable in admin dashboard
        [PROPERTY_EDITABLE => true]
    );

    $head->addMeta('description', $description->toString());
\end{verbatim}

\subsubsection*{Editing Homepage Button Behavior}
Locate the src/components/Home/Home.phtml template and edit the HTML of the page.
To add JavaScript functionality locate home.js in the same directory.
If the script file is created it needs to be created by hand and this line of code added to the template:

\begin{verbatim}
    // getRosource function locates the file
    // relative to the class definition of the $this object
    Javascript::import($this->getResource('home.js'));
\end{verbatim}

Edit the script to the needs and reload the page.
The changes should take place instantaneously.

\subsubsection*{Uploading Files}
Files can be uploaded via the administration dashboard.
The files are limited to pages created in dashboard and can not be used natively in homepage template.
Navigate to the Admin/File System/Files page and drag-and-drop your files on to the grid.
Drag-and-dropping is the only way to upload files to the system.

\subsubsection*{Changing Homepage Image}
Add images to the public/images directory, locate the src/components/Home/Home.phtml template, and create or edit the
desired HTML image element.
The images are accessible on url: <project-domain>/public/images/<image-file>

\subsubsection*{Creating Editor User}
Navigate to the System/Groups and create new group called Editors and fill out the privileges mapping.
Each checkbox corresponds to Resource and Action performed on the resource.
If the checkbox is checked it means that the group has privilege to perform the action on the resource.

Navigate to the System/Users and create new user.
When creating a new user, the Tag needs to be unique for the website.
Add the user into the Admin group which gives them the access to log into the administration dashboard
and the newly created Editor group defines all the actions the user can perform in the dashboard.
